\documentclass[a4paper,12pt]{article}
\usepackage{HomeWorkTemplate}
\usepackage{circuitikz}
\usepackage[shortlabels]{enumitem}
\usepackage{hyperref}
\usepackage{tikz}
\usepackage{amsmath}
\usepackage{amssymb}
\usepackage{tcolorbox}
\usepackage{xepersian}
\settextfont{Persian Modern}
\usetikzlibrary{arrows,automata}
\usetikzlibrary{circuits.logic.US}
\usepackage{changepage}
\newcounter{problemcounter}
\newcounter{subproblemcounter}
\setcounter{problemcounter}{1}
\setcounter{subproblemcounter}{1}
\newcommand{\problem}[1]
{
	\subsection*{
		پرسش
		\arabic{problemcounter} 
		\stepcounter{problemcounter}
		\setcounter{subproblemcounter}{1}
		#1
	}
}
\newcommand{\subproblem}{
	\textbf{\harfi{subproblemcounter})}\stepcounter{subproblemcounter}
}


\begin{document}
\handout
{نظریه زبان‌ها و اتوماتا}
{دکتر شهرام خزایی}
{نیم‌سال دوم ۱۴۰۱-۱۴۰۰}
{}
{}
{زمان: ۲ ساعت و ۴۵ دقیقه}
 {امتحان پایان‌ترم}
 
\textbf{ حتما پیش از پاسخ‌دادن به نکات زیر توجه کنید:}\\
\textcolor{black}{ توجه کنید که برای گرفتن نمره‌ی کامل از پرسش‌ ۱، باید برای ادعاهایتان دلیل کافی بیاورید. با این‌حال، با توجه به محدودیت زمان، به شما اکیدا توصیه می‌شود که اثبات‌هایی مفصل و با جزئیات کامل ارائه نکنید. سعی کنید نکات مهم را به‌طور مختصر خاطرنشان کنید و کلیت طرح اثبات را بیان کنید. در مواردی که برای پاسخ به پرسشی یک الگوریتم، ماشین محاسبه، گرامر، عبارت منظم و ... ارائه می‌کنید، نیاز به اثبات درستی آن‌چه ارائه کرده‌اید نیست؛ اما سعی کنید با توضیحی مختصر، شهود پشت طراحیتان را به مصحح منتقل کنید. همچنین اگر برای پاسخ به قسمتی از پرسش اول، خواستید از لم تزریق استفاده کنید نیز توجه داشته باشید که نیازی نیست درگیر جزئیات حالت‌بندی‌های مختلف بشوید و صرفا بیان کلیت روش، کفایت می‌کند.\\
	در پاسخ به سوال‌های ۲ و ۳، برخلاف پرسش ۱، از شما انتظار می‌رود به‌طور دقیق ادعاهایتان را اثبات کنید. با این همه، اگر نیازی به طراحی ماشین تورینگ داشتید، چه در این سوال و چه در پرسش‌های قبل آن، می‌توانید از توصیف سطح بالا (الگوریتمی) ماشین تورینگ استفاده کنید و نیازی به طراحی سطح پایین (طراحی ساختار ماشین و تابع انتقال آن) نیست!}

\problem{درست و نادرست}
درستی یا نادرستی هر یک از گزاره‌های زیر را با ذکر اجمالی دلیل، تعیین کنید. \textcolor{red}{(هر قسمت ۲ نمره)}\\
\begin{enumerate}
\item 
 اگر در گراف انتقال حالت متناظر با یک $ DFA$، جهت یال‌ها را برعکس کنیم، گراف حاصل نمایش یک $DFA$ است.
 \item
 اگر در گراف انتقال حالت متناظر با یک $ DFA$، جهت یال‌ها را برعکس کنیم، گراف حاصل نمایش یک $NFA$ است.
\item 
 یک زبان منظم، نامتناهی است اگر و فقط اگر گراف انتقال حالت متناظر با $ DFA$ای که آن را می‌پذیرد دارای دور باشد. 
\item 
مساله‌ی متناهی بودن یک زبان منظم، تصمیم‌پذیر است.
\item 
یک ماشین محاسبه‌ی غیرقطعی، ماشینی است که قطعی نیست.
\item 
زبان 
$$L_1=\{1^x**1^y=1^z|x,y,z\in \mathbb{N}~,~z=x^y\}$$
مستقل از متن، ولی غیرمنظم است.
\item 
زبان 
$$L_2=\{ a^ib^j c^k|i,j,k\geq 0~,~ j\neq i~ \text{یا}~ j\neq k \}$$
زبانی مستقل از متن است.
\item 
زبان 
$$L_3=\{w\in \{0,1\}^*| \text{  $w$ کدینگ یک ماشین تورینگ است } \}$$
یک زبان تورینگ-تصمیم‌پذیر نیست.
\item 
مسأله‌ی عضویت رشته‌ی $w$ برای گرامر مستقل از متن داده‌شده، تصمیم‌پذیر است.
\item 
مسأله‌ی تهی بودن زبان یک 
$PDA$
که توصیف آن داده ‌شده است، تصمیم‌ناپذیر است.
\item 
مسأله‌ی برابری زبان یک گرامر مستقل از متن داده شده با 
$\Sigma^*$، 
تورینگ-تشخیص‌پذیر است اما تورینگ-تصمیم‌پذیر نیست. 
\item
اگر هیچ‌یک از زبان‌های 
$L_1,L_2,\dots$
با هیچ 
$DFA$ای 
پذیرفته نشوند و به‌ازای هر 
$i\geq 1$ 
داشته باشیم: 
$L_i\subseteq L_{i+1}$، 
آن‌گاه زبان 
$\cup_{i=1}^{\infty}L_i$ با هیچ 
$DFA$ای پذیرفته نمی‌شود.
\end{enumerate}


\problem{زبان‌های ساده‌ی لا ینحل (؟!	)}
آیا روی الفبای تک‌عضوی 
$\Sigma=\{a\}$، زبان‌هایی وجود دارند که تورینگ-تشخیص‌پذیر نیستند؟\\
ادعای خود را به‌طور کامل اثبات کنید. \textcolor{red}{ (10 نمره)}\\

\problem{مقایسه‌ی سختی}
\begin{enumerate}
	\item 
	دو زبان زیر را در نظر بگیرید:
	$$A=\{\langle M,w\rangle~:~\text{$M$
		یک ماشین تورینگ است و رشته‌ی 
		$w$
		را می‌پذیرد}\}$$
	$$H=\{\langle M,w\rangle~:~\text{$M$
		یک ماشین تورینگ است و روی رشته‌ی 
		$w$
		متوقف می‌شود}\}$$
	با فرض این‌که زبان $A$، تصمیم‌ناپذیر است، تصمیم‌ناپذیری زبان دوم را نشان دهید. \textcolor{red}{ (6 نمره)}
	\item 
	نشان دهید زبان زیر تصمیم‌ناپذیر است:
	$$L_{comp}=\{\langle M_1,M_2\rangle~:~ L(M_1)=L(M_2)^c\}$$
	توجه کنید که در بالا، منظور از 
	$\langle M_1,M_2\rangle$، کدینگ چسبانده‌شده و قابل بازیابی ماشین‌های تورینگ 
	$M_1$
	و 
	$M_2$
	می‌باشد.\\
	شاید برای پاسخ، تصمیم‌ناپذیری زبان 
	$$L_{all}=\{\langle M\rangle~:~L(M)=\Sigma^* \}$$
	بتواند کمکتان کند. \textcolor{red}{ (10 نمره)}\\
\end{enumerate}

\end{document}